\scanpage{43}
\mspage{20}
le groupe qu'ils forment, et par $\mathrm{Autstr}^{\unsure{\circ|0}}(P_M)$ le sous-groupe formé des $u$ qui induisent l'identité sur $\pi_0(P_M)$. On trouve donc une suite exacte canonique
\[
1\longrightarrow\mathrm{Autstr}^\circ(P_M)\longrightarrow\mathrm{Autstr}(P_M)\xrightarrow{\delta} Z'_\Lambda\longrightarrow1,\tag*{(102)}
\]
où l'hom $\delta$ est défini via l'opération d'un $\deleted{\unsure{P_M\ \mathrm{sur}}}u\in\mathrm{Autstr}(P_M)$ sur $\pi_0(P_M)$, qui doit être de la forme $i\mapsto\deleted{\unsure{\lambda_\delta}}\delta(u).i$, pour un $\delta(u)\in Z'_\Lambda$ bien déterminé (puisque $u$ \deleted{conserve}\inserted{respecte} la structure (97)). Considérons maintenant l'hom canonique de suites exactes, \deleted{déduit} de la suite exacte
\[
1\longrightarrow\mathrm{Gl}'_\Lambda(M)^\pm\longrightarrow\mathrm{Gl}_\Lambda(M)'\longrightarrow\Lambda^*/\pm1_\Lambda\unsure{\text{---}|\longrightarrow}1\tag*{(103)}
\]
dans (102), induite par $\dot g\mapsto T_{\dot g}$\,:

\diagram{\small\[\begin{tikzcd}[ampersand replacement=\&,column sep=small,row sep=large]
1\arrow[r]\&\mathrm{Gl}'_\Lambda(M)^\pm\arrow[r]\arrow[d]\&\mathrm{Gl}_\Lambda(M)'\arrow[r]\arrow[d,"{\dot g\mapsto T_{\dot g}}"]\&\Lambda^*/\{\pm1_\Lambda\}\arrow[r]\arrow[d,"1"]\&1\\
1\arrow[r]\&\mathrm{Autstr}^\circ(P_M)\arrow[r]\&\mathrm{Autstr}(P_M)\arrow[r]\&Z'_\Lambda\arrow[r]\&1
\end{tikzcd}\tag*{(104)}\]}

\noindent\begin{tabular}{|p{.91\linewidth}}
\uline{Théorème}\quad L'hom de suites exactes (104) est un iso.
\end{tabular}

L'injectivité de $\dot g\mapsto T_{\dot g}$ est immédiate, elle provient de l'injectivité \deleted{de} $\mathrm{Gl}'_\Lambda(M)^\pm\longrightarrow\mathrm{Aut}(SP_{(M,e_M)})$, déjà vue (cette dernière flèche étant en fait bijective).
% archive sheet number: 42
